\documentclass[11pt]{article}

\usepackage[margin=1in]{geometry}
\usepackage{amsmath,amssymb}
\usepackage{algorithm}
\usepackage{algpseudocode}
\usepackage{booktabs}
\usepackage[colorlinks=true,linkcolor=blue,citecolor=blue,urlcolor=blue]{hyperref}

\newcommand{\N}{\mathbb{N}}
\newcommand{\R}{\mathbb{R}}
\newcommand{\E}{\mathbb{E}}
\newcommand{\Poi}{\mathrm{Poi}}
\newcommand{\Gam}{\mathrm{Gam}}
\newcommand{\KL}{\mathrm{KL}}

\title{Variational Empirical-Bayes Poisson NMF\\
and Effective-Rank Selection}
\author{}
\date{}

\begin{document}

\maketitle

\section{Purpose and main qualification}

Let $Y\in\N^{N\times M}$ be a count matrix.  We begin with a deliberately
large candidate rank $K_{\max}$ and investigate whether factor-specific Gamma
priors on the loading columns can suppress unnecessary factors.  The proposed
model uses
\[
  L_{ik}\mid \alpha_k^0,\beta_k^0
  \overset{\mathrm{ind.}}{\sim}
  \Gam(\alpha_k^0,\beta_k^0),
\]
where the Gamma distribution is parameterized by shape and rate.  A
mean-field approximation is used for $L$, and $(\alpha_k^0,\beta_k^0)$ are
estimated by variational empirical Bayes (VEB).

This procedure keeps $K_{\max}$ fixed and may estimate a smaller
\emph{effective rank}.  It does not put a posterior distribution directly on
$K$.  An unused factor has a VEB fixed-point direction in which its prior mean
$\alpha_k^0/\beta_k^0$ approaches zero.  Nevertheless, unconstrained VEB is
not guaranteed to remove redundant factors: two similar factors may divide
the signal and both retain positive prior means.  Section~\ref{sec:rank}
explains this qualification precisely.


\section{Poisson NMF model}

The observation model is
\begin{align}
  Y_{im}\mid L,F
  &\overset{\mathrm{ind.}}{\sim}
  \Poi\!\left(\sum_{k=1}^{K_{\max}}L_{ik}F_{km}\right),
  \label{eq:observation-model}\\
  L_{ik}\mid\alpha_k^0,\beta_k^0
  &\overset{\mathrm{ind.}}{\sim}
  \Gam(\alpha_k^0,\beta_k^0).
  \label{eq:loading-prior}
\end{align}
We impose
\begin{equation}
  F_{km}\geq 0,
  \qquad
  \sum_{m=1}^M F_{km}=1,
  \qquad k=1,\ldots,K_{\max}.
  \label{eq:F-normalization}
\end{equation}
Under this normalization, $L_{ik}$ is the expected total count contributed by
factor $k$ to observation $i$.  Consequently, the prior mean
$\alpha_k^0/\beta_k^0$ has a directly interpretable scale.

The model can be initialized with vanilla Poisson NMF.  If the initial
factorization is $Y\approx \widetilde L\widetilde F$, define
\begin{equation}
  F_{km}^{(0)}
  =
  \frac{\widetilde F_{km}}
       {\sum_{m'=1}^M\widetilde F_{km'}},
  \qquad
  L_{ik}^{(0)}
  =
  \widetilde L_{ik}\sum_{m=1}^M\widetilde F_{km}.
  \label{eq:nmf-normalization}
\end{equation}
This removes the ordinary NMF scale ambiguity without changing the fitted
mean.


\section{Poisson allocation augmentation}

Introduce latent component-specific counts
\begin{equation}
  Z_{imk}\mid L,F\sim\Poi(L_{ik}F_{km}),
  \qquad
  Y_{im}=\sum_{k=1}^{K_{\max}}Z_{imk}.
  \label{eq:latent-counts}
\end{equation}
Conditional on $Y_{im}$, their variational distribution is multinomial:
\begin{equation}
  q(Z_{im1},\ldots,Z_{imK_{\max}}\mid Y_{im})
  =
  \operatorname{Multinomial}
  \left(Y_{im};\xi_{im1},\ldots,\xi_{imK_{\max}}\right),
  \label{eq:q-Z}
\end{equation}
where $\xi_{imk}\geq0$ and $\sum_k\xi_{imk}=1$.

We use the mean-field family
\begin{equation}
  q(L,Z)
  =
  \left\{
    \prod_{i=1}^N\prod_{k=1}^{K_{\max}}
    \Gam(L_{ik};\alpha_{ik},\beta_{ik})
  \right\}
  \left\{
    \prod_{i=1}^N\prod_{m=1}^M q(Z_{im}\mid Y_{im})
  \right\}.
  \label{eq:mean-field-family}
\end{equation}


\section{Variational lower bound}

Up to the constants $-\sum_{i,m}\log(Y_{im}!)$, the augmented variational
lower bound is
\begin{align}
  \mathcal L
  ={}&
  \sum_{i=1}^N\sum_{m=1}^M\sum_{k=1}^{K_{\max}}
  Y_{im}\xi_{imk}
  \left\{
    \psi(\alpha_{ik})-\log\beta_{ik}
    +\log F_{km}-\log\xi_{imk}
  \right\}
  \nonumber\\
  &-
  \sum_{i=1}^N\sum_{k=1}^{K_{\max}}
  \frac{\alpha_{ik}}{\beta_{ik}}
  \sum_{m=1}^M F_{km}
  \nonumber\\
  &-
  \sum_{i=1}^N\sum_{k=1}^{K_{\max}}
  \KL\!\left\{
    \Gam(\alpha_{ik},\beta_{ik})
    \,\middle\|\,
    \Gam(\alpha_k^0,\beta_k^0)
  \right\}.
  \label{eq:elbo}
\end{align}
We use the conventions $0\log0=0$ and $\log F_{km}=-\infty$ when
$F_{km}=0$.  In computation, the log is protected by a small numerical
constant, but the factor update itself may preserve exact zeros if desired.

For Gamma distributions parameterized by shape and rate,
\begin{align}
  &\KL\!\left\{
    \Gam(\alpha,\beta)
    \,\middle\|\,
    \Gam(\alpha^0,\beta^0)
  \right\}
  \nonumber\\
  &\quad=
  (\alpha-\alpha^0)\psi(\alpha)
  +\alpha^0\log\frac{\beta}{\beta^0}
  -\log\Gamma(\alpha)+\log\Gamma(\alpha^0)
  -\alpha+\frac{\alpha\beta^0}{\beta}.
  \label{eq:gamma-kl}
\end{align}


\section{Coordinate updates}

\subsection{Latent count allocations}

For every $i,m,k$,
\begin{equation}
  \xi_{imk}
  =
  \frac{
    F_{km}\exp\{\psi(\alpha_{ik})-\log\beta_{ik}\}
  }{
    \sum_{h=1}^{K_{\max}}
    F_{hm}\exp\{\psi(\alpha_{ih})-\log\beta_{ih}\}
  }.
  \label{eq:xi-update}
\end{equation}
Only entries with $Y_{im}>0$ contribute to subsequent sufficient statistics,
so the implementation need not store $\xi_{imk}$ for zero observations.

\subsection{Loading posterior}

Conjugacy gives
\begin{align}
  \alpha_{ik}
  &\leftarrow
  \alpha_k^0+
  \sum_{m=1}^M Y_{im}\xi_{imk},
  \label{eq:alpha-update}\\
  \beta_{ik}
  &\leftarrow
  \beta_k^0+
  \sum_{m=1}^M F_{km}.
  \label{eq:beta-update}
\end{align}
Because of Eq.~\eqref{eq:F-normalization},
\begin{equation}
  \beta_{ik}=\beta_k^0+1.
  \label{eq:beta-normalized}
\end{equation}

\subsection{Factor dictionary}

Let
\begin{equation}
  C_{km}
  =
  \sum_{i=1}^N Y_{im}\xi_{imk}.
  \label{eq:factor-count}
\end{equation}
Maximizing the lower bound over a row-normalized $F$ gives
\begin{equation}
  F_{km}
  \leftarrow
  \frac{C_{km}}{\sum_{m'=1}^M C_{km'}}.
  \label{eq:F-update}
\end{equation}
If $\sum_m C_{km}$ is numerically zero, the factor is provisionally inactive;
its previous row may be retained until the final pruning decision.  Damping
the update can improve stability when factors are nearly redundant.

\subsection{Variational empirical-Bayes update}

For each candidate factor $k$, define the posterior averages
\begin{align}
  \overline{\lambda}_k
  &=
  \frac{1}{N}\sum_{i=1}^N
  \frac{\alpha_{ik}}{\beta_{ik}},
  \label{eq:mean-loading}\\
  \overline{\log\lambda}_k
  &=
  \frac{1}{N}\sum_{i=1}^N
  \left\{\psi(\alpha_{ik})-\log\beta_{ik}\right\}.
  \label{eq:mean-log-loading}
\end{align}
Maximizing
\[
  \sum_{i=1}^N
  \E_q\!\left[
    \log\Gam(L_{ik};\alpha_k^0,\beta_k^0)
  \right]
\]
over $(\alpha_k^0,\beta_k^0)$ gives
\begin{align}
  \beta_k^0
  &\leftarrow
  \frac{\alpha_k^0}{\overline{\lambda}_k},
  \label{eq:prior-rate-update}\\
  \log\alpha_k^0-\psi(\alpha_k^0)
  &=
  \log\overline{\lambda}_k-
  \overline{\log\lambda}_k.
  \label{eq:prior-shape-equation}
\end{align}
Equation~\eqref{eq:prior-shape-equation} can be solved by Newton's method.  If
\[
  g(\alpha)
  =
  \log\alpha-\psi(\alpha)
  -\log\overline{\lambda}_k
  +\overline{\log\lambda}_k,
\]
then
\begin{equation}
  \alpha
  \leftarrow
  \alpha-
  \frac{g(\alpha)}{1/\alpha-\psi_1(\alpha)},
  \label{eq:newton-update}
\end{equation}
where $\psi_1$ is the trigamma function.  Numerical lower and upper bounds on
$\alpha_k^0$ and damping of both hyperparameter updates are recommended.


\section{Inference algorithm}

\begin{algorithm}[H]
\caption{VEB Poisson NMF initialized by vanilla NMF}
\label{alg:veb-nmf}
\small
\begin{algorithmic}[1]
\Require Count matrix $Y$, candidate rank $K_{\max}$, inner iterations $R$,
outer iterations $T$, convergence tolerance $\delta$, and numerical constant
$\epsilon$.
\State Fit rank-$K_{\max}$ vanilla Poisson NMF to obtain
$(\widetilde L,\widetilde F)$.
\State Normalize the factor rows and rescale the loading columns using
Eq.~\eqref{eq:nmf-normalization}; call the results $(L^{(0)},F^{(0)})$.
\State Initialize $\xi_{imk}\propto L_{ik}^{(0)}F_{km}^{(0)}$.
\State Initialize positive $(\alpha_k^0,\beta_k^0)$ from the columns of
$L^{(0)}$, and initialize $(\alpha_{ik},\beta_{ik})$ using
Eqs.~\eqref{eq:alpha-update}--\eqref{eq:beta-update}.
\For{$t=1,\ldots,T$}
  \For{$r=1,\ldots,R$}
    \State Update $\xi$ using Eq.~\eqref{eq:xi-update}, retaining only the
    sufficient statistics needed for nonzero $Y_{im}$.
    \State Update $(\alpha_{ik},\beta_{ik})$ using
    Eqs.~\eqref{eq:alpha-update}--\eqref{eq:beta-update}.
  \EndFor
  \State Update $F$ using Eq.~\eqref{eq:F-update}; optionally damp and then
  row-normalize.
  \State Update $(\alpha_k^0,\beta_k^0)$ independently for every $k$ using
  Eqs.~\eqref{eq:prior-rate-update}--\eqref{eq:newton-update}; optionally damp.
  \State Evaluate the lower bound in Eq.~\eqref{eq:elbo}.
  \If{the relative lower-bound improvement is less than $\delta$}
    \State \textbf{break}
  \EndIf
\EndFor
\State Recompute $\xi$ and $q(L)$ using the final $F$ and prior parameters.
\State Calculate the relevance summaries in Eqs.~\eqref{eq:relevance-mean}
and~\eqref{eq:allocated-share}, and prune only after convergence.
\State \Return $F$, $(\alpha_{ik},\beta_{ik})$,
$(\alpha_k^0,\beta_k^0)$, the lower-bound history, and
$\widehat K_{\mathrm{eff}}$.
\end{algorithmic}
\end{algorithm}


\section{Can this procedure learn the rank?}
\label{sec:rank}

The natural VEB relevance summary is
\begin{equation}
  r_k
  =
  \frac{\alpha_k^0}{\beta_k^0}.
  \label{eq:relevance-mean}
\end{equation}
At an interior VEB optimum, $r_k=\overline\lambda_k$.  A complementary
data-allocation summary is
\begin{equation}
  p_k
  =
  \frac{
    \sum_{i=1}^N\sum_{m=1}^M Y_{im}\xi_{imk}
  }{
    \sum_{i=1}^N\sum_{m=1}^M Y_{im}
  }.
  \label{eq:allocated-share}
\end{equation}
For an exploratory implementation, one may report
\begin{equation}
  \widehat K_{\mathrm{eff}}(\tau)
  =
  \#\left\{
    k:
    \frac{r_k}{\sum_{h=1}^{K_{\max}}r_h}\geq\tau
    \ \text{and}\ 
    p_k\geq\tau
  \right\}.
  \label{eq:effective-rank}
\end{equation}
This is a thresholded effective rank, not a posterior estimate of $K$.

There is a real shrinkage mechanism.  Suppose factor $k$ receives no counts,
so that $\sum_mY_{im}\xi_{imk}=0$ for every $i$.  Then
\[
  \alpha_{ik}=\alpha_k^0,
  \qquad
  \beta_{ik}=\beta_k^0+1.
\]
The next undamped VEB update leaves the shape unchanged and maps
\[
  \beta_k^0\leftarrow\beta_k^0+1.
\]
Repeated updates therefore drive $r_k=\alpha_k^0/\beta_k^0$ toward zero.  In
this sense, an empty component can be automatically suppressed.

The difficulty is that an overfitted component need not remain empty.  The
factor update may adapt it to sampling noise, after which it receives counts
and escapes shrinkage.  More fundamentally, suppose two factor rows are
identical, $F_{1\cdot}=F_{2\cdot}$, and their loading priors have a common rate:
\[
  L_{i1}\sim\Gam(\alpha_1^0,\beta^0),
  \qquad
  L_{i2}\sim\Gam(\alpha_2^0,\beta^0).
\]
Then
\begin{equation}
  L_{i1}+L_{i2}
  \sim
  \Gam(\alpha_1^0+\alpha_2^0,\beta^0).
  \label{eq:gamma-splitting}
\end{equation}
The likelihood cannot distinguish this two-factor representation from a
single identical factor having shape $\alpha_1^0+\alpha_2^0$.  Consequently,
VEB may have both a pruned solution and a signal-splitting solution, with the
selected solution depending on initialization and optimization.


\section{Relationship to conventional ARD}

In the ARD-NMF model of \cite{tan2013ard}, a common relevance scale ties an
entire dictionary component to its corresponding loading vector.  Optimizing
that scale yields a group-sparsity penalty, and a component is pruned when both
sides approach zero.  The VEB model above differs in two ways: $F$ is
constrained to have unit row sums, and $(\alpha_k^0,\beta_k^0)$ are fitted to
the loading distribution without a sparsity-inducing hyperprior.  Learning a
prior is therefore not, by itself, equivalent to ARD.

A closer conjugate ARD extension would fix a common loading shape $a_L$ and
introduce a factor-specific rate,
\begin{align}
  L_{ik}\mid\tau_k
  &\sim\Gam(a_L,\tau_k),\\
  \tau_k
  &\sim\Gam(c_0,d_0).
\end{align}
Under a mean-field approximation,
\begin{equation}
  q(\tau_k)
  =
  \Gam\!\left(
    c_0+Na_L,
    d_0+\sum_{i=1}^N\E_q[L_{ik}]
  \right).
  \label{eq:hierarchical-ard-update}
\end{equation}
A factor with little total loading obtains a large posterior rate
$\E_q[\tau_k]$, which then shrinks every $L_{ik}$ in that column.  A global
spike-and-slab indicator for each loading column would provide exact component
exclusion, at the cost of a more involved variational update.

For the first experiment, the unconstrained VEB algorithm should be
implemented exactly as Algorithm~\ref{alg:veb-nmf}, because its successes and
failures are scientifically informative.  It should then be compared with
(i) fixed $\alpha_k^0$ and learned $\beta_k^0$, and (ii) the hierarchical
relevance-rate model in Eq.~\eqref{eq:hierarchical-ard-update}.


\section{Proposed validation experiments}

For data generated with true rank one, fit
$K_{\max}\in\{1,2,5,10\}$ and use multiple random initializations.  Record:
\begin{enumerate}
  \item $r_k=\alpha_k^0/\beta_k^0$ and the allocated count share $p_k$;
  \item the thresholded effective rank over a range of $\tau$;
  \item training and held-out Poisson deviance;
  \item pairwise cosine similarities among the learned rows of $F$;
  \item whether redundant factors are pruned or split the true signal;
  \item sensitivity to learning both Gamma parameters versus fixing the shape;
  \item sensitivity to initialization and update damping.
\end{enumerate}
Held-out prediction is important because a redundant factor may survive by
fitting sampling noise in the same data used to learn $F$ and its prior.


\begin{thebibliography}{1}

\bibitem{tan2013ard}
V.~Y.~F. Tan and C. F\'evotte.
\newblock Automatic relevance determination in nonnegative matrix
factorization with the $\beta$-divergence.
\newblock \emph{IEEE Transactions on Pattern Analysis and Machine
Intelligence}, 35(7):1592--1605, 2013.
\newblock \url{https://arxiv.org/abs/1111.6085}.

\end{thebibliography}

\end{document}
